\begin{abstract}
Decentralized replicated systems face the fundamental challenge of consistently
merging concurrent changes to shared data without relying on a central
coordinating authority. The CAP theorem \cite{brewer2000} shows that in
partitioned networks, strong consistency and high availability cannot be
guaranteed simultaneously \cite{gilbert2002}, which is why modern distributed
systems frequently adopt Eventual Consistency as a core design goal
\cite{shapiro2011comprehensive}. In this context, two competing approaches to
automatic conflict resolution have emerged: Conflict-free Replicated Data Types
(CRDTs) \cite{shapiro2011comprehensive} and Operational Transformation (OT)
\cite{ellis1989}.

This paper analyzes both approaches. It has the focus on their theoretical
foundations, practical implementation complexity, and suitability for
decentralized environments with high concurrency. Both state-based and
operation-based CRDTs are explained and contrasted with established OT control
algorithms \cite{shapiro2011crdt}. The evaluation is conducted based on the
criteria of latency, convergence behavior, scalability, and memory consumption.

The results indicate that CRDTs, because of their mathematically guaranteed
convergence properties, work best for serverless peer-to-peer architectures
\cite{kleppmann2019}. On the other hand, OT-based approaches demonstrate their
strengths primarily in server-coordinated collaborative systems
\cite{ellis1989}. Finally, hybrid approaches such as Yjs \cite{nicolaescu2015}
are discussed, which combine properties of both paradigms to compensate for
their respective weaknesses.
\end{abstract}